\section*{第 \ref{ch:g12:adaptive-immunity} 章——适应性免疫与疫苗接种}

\begin{solution}{exo:g12:adaptive-immunity:1}
\omterm{def:g12:adaptive-immunity:antigen}{抗原}是\omterm{def:g12:innate-immunity:immunity}{适应性免疫}系统能够识别的任何分子；
\omterm{def:g12:adaptive-immunity:antigen}{淋巴细胞}是携带针对某一种\omterm{def:g12:adaptive-immunity:antigen}{抗原}受体的白\omterm{def:g10:cells-common-unit:cell}{细胞}。其
专一性来自其受体，该受体在成熟过程中随机装配，
且其所有子代\omterm{def:g10:cells-common-unit:cell}{细胞}的受体相同。
\end{solution}

\begin{solution}{exo:g12:adaptive-immunity:2}
\omterm{def:g12:adaptive-immunity:antigen}{抗原}与少数受体匹配的\omterm{def:g12:adaptive-immunity:antigen}{淋巴细胞}结合；这些\omterm{def:g10:cells-common-unit:cell}{细胞}
在数日内增殖形成克隆；该克隆分化为
执行战斗任务的效应\omterm{def:g10:cells-common-unit:cell}{细胞}；少数\omterm{def:g10:cells-common-unit:cell}{细胞}成为\omterm{prop:g12:adaptive-immunity:selection}{记忆细胞}，
长期存活。
\end{solution}

\begin{solution}{exo:g12:adaptive-immunity:3}
由浆\omterm{def:g10:cells-common-unit:cell}{细胞}分泌的 Y 形\omterm{def:g11:gene-expression:protein}{蛋白质}，是 B
\omterm{def:g10:cells-common-unit:cell}{细胞}受体的可溶形式；每条臂都能结合\omterm{def:g12:adaptive-immunity:antigen}{抗原}。附着于细菌表面时，它形成
一层被\omterm{def:g12:innate-immunity:phagocytosis}{吞噬细胞}抓住的包被，并将细菌交联成复合体，
随后被吞噬。
\end{solution}

\begin{solution}{exo:g12:adaptive-immunity:4}
辅助性 T \omterm{def:g10:cells-common-unit:cell}{细胞}在\omterm{prop:g12:innate-immunity:handover}{树突状细胞}激活下分泌\omterm{def:g10:cells-common-unit:cell}{细胞}因子，
这些\omterm{def:g10:cells-common-unit:cell}{细胞}因子可授权并放大 B \omterm{def:g10:cells-common-unit:cell}{细胞}和其他 T \omterm{def:g10:cells-common-unit:cell}{细胞}的应答。
\omterm{def:g10:cells-common-unit:cell}{细胞}毒性 T \omterm{def:g10:cells-common-unit:cell}{细胞}杀死表面展示病毒或
\omterm{def:g11:cancer:cancer}{肿瘤}蛋白片段的\omterm{def:g10:cells-common-unit:cell}{细胞}。
\end{solution}

\begin{solution}{exo:g12:adaptive-immunity:5}
病原体\omterm{def:g12:adaptive-immunity:antigen}{抗原}的无害形式（减毒、灭活、
\omterm{def:g11:gene-expression:protein}{蛋白质}、失活毒素或其遗传指令）并加入佐剂。
它引发初次应答并留下\omterm{prop:g12:adaptive-immunity:selection}{记忆细胞}，使得
真正的病原体入侵时可引发再次应答。
\end{solution}

\begin{solution}{exo:g12:adaptive-immunity:6}
初次应答：延迟约一周，峰值较低（约为
再次的十分之一），两至三周时出现，一个月内回落。
再次应答：延迟两至三天，十天内出现的峰值
高十倍，并持续数月。
\end{solution}

\begin{solution}{exo:g12:adaptive-immunity:7}
记忆具有专一性：由 A 的首次剂量留下的\omterm{prop:g12:adaptive-immunity:selection}{记忆细胞}
只应答 A；初次遇到 B 时，则产生初次应答。
\end{solution}

\begin{solution}{exo:g12:adaptive-immunity:8}
\omterm{prop:g12:adaptive-immunity:antibodies}{抗体}是体液中的\omterm{def:g11:gene-expression:protein}{蛋白质}，不能进入\omterm{def:g10:cells-common-unit:cell}{细胞}。
被感染的\omterm{def:g10:cells-common-unit:cell}{细胞}在表面展示病毒片段，
\omterm{def:g10:cells-common-unit:cell}{细胞}毒性 T \omterm{def:g10:cells-common-unit:cell}{细胞}识别这些片段并杀死该\omterm{def:g10:cells-common-unit:cell}{细胞}。
\end{solution}

\begin{solution}{exo:g12:adaptive-immunity:9}
$1 - 1/4 = 75\%$; $1 - 1/18 \approx 94\%$.
\end{solution}

\begin{solution}{exo:g12:adaptive-immunity:10}
阈值以上时，每例病例感染的人数少于一人，
病毒无法持续传播，很少传到婴儿；免疫的
成人构成了围绕婴儿的屏障。
\end{solution}

\begin{solution}{exo:g12:adaptive-immunity:11}
每微升约 750、520 和 180 个；\omterm{prop:g12:adaptive-immunity:hiv}{艾滋病}约在第 9 年，
即计数跌破 200 时开始。
\end{solution}

\begin{solution}{exo:g12:adaptive-immunity:12}
辅助性 T \omterm{def:g10:cells-common-unit:cell}{细胞}分泌的\omterm{def:g10:cells-common-unit:cell}{细胞}因子授权 B \omterm{def:g10:cells-common-unit:cell}{细胞}
分化为浆\omterm{def:g10:cells-common-unit:cell}{细胞}，并授权\omterm{def:g10:cells-common-unit:cell}{细胞}毒性 T \omterm{def:g10:cells-common-unit:cell}{细胞}增殖；两个分支都依赖
同一协调者，因此去除它会使两者同时失效。
\end{solution}

\begin{solution}{exo:g12:adaptive-immunity:13}
\omterm{prop:g12:adaptive-immunity:antibodies}{抗体}是针对病毒当时的形态选择的；病毒
表面蛋白\omterm{def:g11:mutations:mutation}{突变}的速度快于新克隆被选择的速度，
因此每次应答都会被超越。此外，病毒还会隐藏并增殖在
正是负责协调应答的辅助性 T \omterm{def:g10:cells-common-unit:cell}{细胞}之中。
\end{solution}

\begin{solution}{exo:g12:adaptive-immunity:14}
\omterm{prop:g12:adaptive-immunity:selection}{记忆细胞}识别\omterm{def:g12:adaptive-immunity:vaccine}{疫苗}所呈递的表面蛋白。
流感病毒的表面蛋白每年都发生\omterm{def:g11:mutations:mutation}{突变}，形成旧记忆无法结合的形式；
麻疹病毒的表面蛋白几乎不变，因此 1970 年获得的记忆
仍然适用于今天的病毒。
\end{solution}

\begin{solution}{exo:g12:adaptive-immunity:15}
变化的是\omterm{def:g12:adaptive-immunity:antigen}{淋巴细胞}的受体，它们随机产生；
起选择作用的是\omterm{def:g12:adaptive-immunity:antigen}{抗原}，它只使匹配的克隆增殖；被遗传的是
受体，通过克隆的后代遗传，包括
\omterm{prop:g12:adaptive-immunity:selection}{记忆细胞}。时间尺度：应答以天计，记忆以一生计，
而\omterm{def:g12:selection-drift-speciation:population}{种群}以世代计。这一类比止于
传递环节：被选择的克隆随个体死亡，不会
传递给下一代。
\end{solution}

\begin{solution}{pb:g12:adaptive-immunity:1}
\textbf{1.} 约一周内无应答，两至三周时出现一个不高的峰值，
之后数周内逐渐下降。

\textbf{2.} 两至三天内即出现\omterm{prop:g12:adaptive-immunity:antibodies}{抗体}，峰值高十倍，
持续数月：首次剂量留下的\omterm{prop:g12:adaptive-immunity:selection}{记忆细胞}
已形成较大的克隆，因此增殖从数千个
\omterm{def:g10:cells-common-unit:cell}{细胞}开始，而不是从少数几个开始。

\textbf{3.} 活病毒会在\omterm{def:g10:cells-common-unit:cell}{细胞}内短暂增殖，因此
被感染\omterm{def:g10:cells-common-unit:cell}{细胞}会展示其片段，\omterm{def:g10:cells-common-unit:cell}{细胞}毒性 T \omterm{def:g10:cells-common-unit:cell}{细胞}
与\omterm{prop:g12:adaptive-immunity:antibodies}{抗体}一同被选择；灭活病毒则主要激活
\omterm{prop:g12:adaptive-immunity:antibodies}{抗体}分支。

\textbf{4.} 已存在的\omterm{prop:g12:adaptive-immunity:antibodies}{抗体}可阻挡大部分病毒；
记忆 B \omterm{def:g10:cells-common-unit:cell}{细胞}与 T \omterm{def:g10:cells-common-unit:cell}{细胞}在两天内增殖；被感染的\omterm{def:g10:cells-common-unit:cell}{细胞}
在病毒扩散之前即被清除。感染在出现
症状之前即被终止。

\textbf{5.} 病毒无声增殖三天，随后出现发热、皮疹和
病症，初次应答在一周内逐渐形成；约第十天恢复，
并获得\omterm{prop:g12:adaptive-immunity:selection}{记忆细胞}与终身免疫——前提是他
挺过并发症。

\textbf{6.} $1 - 1/15 \approx 93\%$.

\textbf{7.} $15 \times 0.10 = 1.5$：每例病例平均感染的人数超过一人，
因此疾病可以扩散。未受保护。

\textbf{8.} $10\% \times \num{50000} = 5000$ 人未获免疫，其中
1500 人为婴儿。

\textbf{9.} 易感者 2500 人；$15 \times 0.05 = 0.75$：不足一人，
疾病随之消亡。

\textbf{10.} 婴儿无法接种\omterm{def:g12:adaptive-immunity:vaccine}{疫苗}，因此他们唯一的保护在于
周围的成人已获免疫，使病毒没有途径
传到他们身上。

\textbf{11.} 第 3 代：$1.5^3 \approx 3.4$ 例病例；第 6 代：
$1.5^6 \approx 11$.

\textbf{12.} $1 + 1.5 + 2.25 + 3.4 + 5.1 + 7.6 + 11.4 \approx 32$ 例病例。

\textbf{13.} 约 30\% 的 32 例病例中，约有 10 名婴儿。

\textbf{14.} 每例病例使一名易感者变为免疫者，
因此易感者比例下降，每例病例所感染的人数也随之减少；
若易感者仍有 5000 人，疫情只有在
数百例之后才会减缓，历时数月，除非进行\omterm{def:g12:adaptive-immunity:vaccine}{疫苗}接种干预。

\textbf{15.} 覆盖度为 95\% 时：$0.75^3 \approx 0.4$ 与 $0.75^6 \approx 0.18$ 例病例；
合计约 $1 + 0.75 + 0.56 + \dots \approx 4$ 例病例。病例数
约少十倍，疫情自行熄灭。

\textbf{16.} 8 年间共 800 例，即每年每微升 100 例，
约每月 8 例。

\textbf{17.} B \omterm{def:g10:cells-common-unit:cell}{细胞}需要辅助性 T \omterm{def:g10:cells-common-unit:cell}{细胞}分泌的\omterm{def:g10:cells-common-unit:cell}{细胞}因子才能分化为浆\omterm{def:g10:cells-common-unit:cell}{细胞}，
\omterm{def:g10:cells-common-unit:cell}{细胞}毒性 T \omterm{def:g10:cells-common-unit:cell}{细胞}也需要它们才能增殖；辅助性\omterm{def:g10:cells-common-unit:cell}{细胞}太少时，
两条分支都无法获得授权，尽管相应的\omterm{def:g10:cells-common-unit:cell}{细胞}已经存在。

\textbf{18.} 是的：\omterm{prop:g12:adaptive-immunity:selection}{记忆细胞}在扩增为二次应答时，仍需要辅助性 T \omterm{def:g10:cells-common-unit:cell}{细胞}的\omterm{def:g10:cells-common-unit:cell}{细胞}因子；
缺少协调，\omterm{prop:g12:adaptive-immunity:selection}{记忆细胞}便无法
发挥作用；麻疹对健康成人无害，却可能致命。

\textbf{19.} 辅助性 T \omterm{def:g10:cells-common-unit:cell}{细胞}群体可由存活\omterm{def:g10:cells-common-unit:cell}{细胞}和骨髓
重新补充，其他分支的\omterm{prop:g12:adaptive-immunity:selection}{记忆细胞}也从未
被破坏，协调功能一旦恢复，\omterm{def:g10:cells-common-unit:cell}{细胞}库便完好无损。
病毒潜伏在隐蔽的\omterm{def:g10:cells-common-unit:cell}{细胞}中，因此停药后病毒会重新增殖。

\textbf{20.} 覆盖度 93\%；在 90\% 覆盖度下，一名旅行者约引发 32 例病例；
辅助性 T \omterm{def:g10:cells-common-unit:cell}{细胞}。
\end{solution}
